\section{Careers Week \& Half Term}

\subsection*{Overview:}
During this extended half term break I continued development on my web application by building the graph view from my previous prototype into this one and then built that graph view into a simple dashboard which 

\subsection*{Achievements, Thoughts and Ideas:}
My main goal for this couple of weeks was to get a fully functional graph view that enable users to load their notes from the server and have them shown to them in a visually appealing way. 

Initially, I'd hoped this would be as simple as copy-pasting the code from my latest prototype into this project unfortunately this didn't quite turn out to be the case. After importing the code I was instantly met by a large list of errors ranging from compilation errors (the browser didn’t understand the code being presented) to missing references errors (information that the program should have been supplied with was missing altogether), this was certainly disheartening. After reading through the error list in more detail I found it was due to the Vue Js system I was now using and how it didn’t interact very well with my current method of referring libraries.

Whilst I could have likely fixed the errors with some 'band aid' solutions I felt it was best to go back to basics and rebuild the graph layout section by section in a new optimised way. Whilst it was frustrating repeating work I had already done I noticed many areas of my program that could be upgraded in minor ways to boost performance (speed of functions), reliability, features and more. From purely this restructuring my code feels "around 70% faster" - according to my father who is at this point very familiar with me referring changes to him for his opinion

After ripping the code out and starting again I soon recreated the functionality I previously had and can work on new designs. I replaced my prior stack of buttons with a much cleaner toolbar that shows the user quick commands with both visual icons and descriptive text. This arrangement makes finding functions much easier for the user, and to cater for mobile / smaller screen devices the toolbar automatically collapses into dropdown menus, giving the same functionality but in a much reduced area.

Following on from the graph creation I wanted to put the modular component within a user facing screen. I quickly made a Dashboard which presents the user with an area for their recent notes, shortcuts (such as a space to store daily notetaking ideas or a to do list) and the all important graph of their 'second brain'

Due to the time spent fixing the code issues I faced, I didn't quite achieve  all of my goals however I don’t feel this time has been wasted in the slightest as my understanding of the fundamental principles behind web development have hugely expanded and I will certainly find cost/time savings in the coming weeks as I produce other elements in more efficient ways. Referring back to my plan I can see that I'm still ahead of schedule and thanks to having contingency time considered, despite these set backs I am still making good progress

\subsection*{What went well?}
Despite setback I now have an excellent foundation which I can build upon
\subsection*{What did I learn?}
The values of 'starting afresh' when faced with technical setbacks that seem impossibly difficult

\subsection*{Plans for next week:}
Next week I hope to build out the note-taking portion of my webapp to allow the user to record their ideas
